% ═════════════════════════════════════════════════════════════════════════════
\section{Binary Linear Codes}
% ═════════════════════════════════════════════════════════════════════════════

A \defn{binary linear code} $C$ is a subspace of $\mathbb{F}_2^n$. It may be
specified as the nullspace of a \defn{parity check matrix} $A$:
\[
    C = \ker(A) = \{\mathbf{x} \in \mathbb{F}_2^n : A\mathbf{x} = \mathbf{0}\}.
\]
The \defn{syndrome} of a received word $\mathbf{r}$ is $A\mathbf{r}$.
If the syndrome is $\mathbf{0}$, then $\mathbf{r} \in C$; otherwise an error
has occurred.

% ═════════════════════════════════════════════════════════════════════════════
\section{The Code and Its Basis}
% ═════════════════════════════════════════════════════════════════════════════

Let $C = \ker(A)$ where
\[
    A =
    \begin{bmatrix}
        1 & 0 & 1 & 0 & 0 & 1 \\
        1 & 1 & 1 & 0 & 1 & 0 \\
        1 & 1 & 0 & 1 & 0 & 0 \\
        0 & 0 & 1 & 1 & 1 & 0
    \end{bmatrix}.
\]

% CORRECTION: The original rref computation was incorrect. The matrix A has
% rank 3 (since R2 = R3 + R4 mod 2), giving a code of dimension 6-3 = 3,
% not dimension 2 as implied by the original (erroneous) rref of rank 4.
% The basis below is corrected accordingly.

\begin{note}
One can verify that $R_2 = R_3 + R_4 \pmod{2}$, so the rows of $A$ are
linearly dependent over $\mathbb{F}_2$ and $\mathrm{rank}(A) = 3$.
Hence $\dim(C) = 6 - 3 = 3$.
\end{note}

\subsection{Row reduction}

Reducing $A$ over $\mathbb{F}_2$ via elementary row operations:
\[
\begin{bmatrix}
    1 & 0 & 1 & 0 & 0 & 1 \\
    1 & 1 & 1 & 0 & 1 & 0 \\
    1 & 1 & 0 & 1 & 0 & 0 \\
    0 & 0 & 1 & 1 & 1 & 0
\end{bmatrix}
\xrightarrow{R_2 \leftarrow R_2+R_1,\; R_3 \leftarrow R_3+R_1}
\begin{bmatrix}
    1 & 0 & 1 & 0 & 0 & 1 \\
    0 & 1 & 0 & 0 & 1 & 1 \\
    0 & 1 & 1 & 1 & 0 & 1 \\
    0 & 0 & 1 & 1 & 1 & 0
\end{bmatrix}
\xrightarrow{R_3 \leftarrow R_3+R_2}
\begin{bmatrix}
    1 & 0 & 1 & 0 & 0 & 1 \\
    0 & 1 & 0 & 0 & 1 & 1 \\
    0 & 0 & 1 & 1 & 1 & 0 \\
    0 & 0 & 1 & 1 & 1 & 0
\end{bmatrix}
\]
\[
\xrightarrow{R_4 \leftarrow R_4+R_3}
\begin{bmatrix}
    1 & 0 & 1 & 0 & 0 & 1 \\
    0 & 1 & 0 & 0 & 1 & 1 \\
    0 & 0 & 1 & 1 & 1 & 0 \\
    0 & 0 & 0 & 0 & 0 & 0
\end{bmatrix}
\xrightarrow{R_1 \leftarrow R_1+R_3}
\underbrace{\begin{bmatrix}
    1 & 0 & 0 & 1 & 1 & 1 \\
    0 & 1 & 0 & 0 & 1 & 1 \\
    0 & 0 & 1 & 1 & 1 & 0 \\
    0 & 0 & 0 & 0 & 0 & 0
\end{bmatrix}}_{\mathrm{rref}(A)}
\]

Pivot variables: $x_1, x_2, x_3$. Free variables: $x_4, x_5, x_6$.

\subsection{Basis for \texorpdfstring{$C$}{C}}

Setting $(x_4, x_5, x_6) = (s, t, u)$ and back-substituting:
\begin{align*}
    x_3 &= x_4 + x_5 = s + t, \\
    x_2 &= x_5 + x_6 = t + u, \\
    x_1 &= x_4 + x_5 + x_6 = s + t + u.
\end{align*}

The general codeword is
\[
    \mathbf{x}
    = s\underbrace{\begin{pmatrix}1\\0\\1\\1\\0\\0\end{pmatrix}}_{\mathbf{v}_1}
    + t\underbrace{\begin{pmatrix}1\\1\\1\\0\\1\\0\end{pmatrix}}_{\mathbf{v}_2}
    + u\underbrace{\begin{pmatrix}1\\1\\0\\0\\0\\1\end{pmatrix}}_{\mathbf{v}_3},
    \qquad s, t, u \in \mathbb{F}_2.
\]

\[
    \text{Basis of } C = \left\{
    \begin{pmatrix}1\\0\\1\\1\\0\\0\end{pmatrix},\;
    \begin{pmatrix}1\\1\\1\\0\\1\\0\end{pmatrix},\;
    \begin{pmatrix}1\\1\\0\\0\\0\\1\end{pmatrix}
    \right\}.
\]

% ═════════════════════════════════════════════════════════════════════════════
\section{Syndrome Computation}
% ═════════════════════════════════════════════════════════════════════════════

\begin{example}
Determine which of the received words
\[
    \mathbf{x} = (1,0,1,0,1,0)^{\top}, \qquad \mathbf{y} = (1,1,1,0,1,0)^{\top}
\]
is a valid codeword of $C$.
\end{example}

\begin{proof}[Solution]
Compute the syndrome $A\mathbf{r} \pmod{2}$ for each received word.

\textbf{Syndrome of $\mathbf{x}$:}
\[
    A\mathbf{x} =
    \begin{pmatrix}
        1{+}0{+}1{+}0{+}0{+}0\\
        1{+}0{+}1{+}0{+}1{+}0\\
        1{+}0{+}0{+}0{+}0{+}0\\
        0{+}0{+}1{+}0{+}1{+}0
    \end{pmatrix}
    \equiv
    \begin{pmatrix}0\\1\\1\\0\end{pmatrix} \pmod{2}.
\]
Since $A\mathbf{x} \neq \mathbf{0}$, $\mathbf{x}$ is \textbf{not} a valid codeword.

\medskip
\textbf{Syndrome of $\mathbf{y}$:}
\[
    A\mathbf{y} =
    \begin{pmatrix}
        1{+}0{+}1{+}0{+}0{+}0\\
        1{+}1{+}1{+}0{+}1{+}0\\
        1{+}1{+}0{+}0{+}0{+}0\\
        0{+}0{+}1{+}0{+}1{+}0
    \end{pmatrix}
    \equiv
    \begin{pmatrix}0\\0\\0\\0\end{pmatrix} \pmod{2}.
\]
Since $A\mathbf{y} = \mathbf{0}$, $\mathbf{y}$ is a \textbf{valid} codeword.
(Indeed, $\mathbf{y} = \mathbf{v}_2$, one of the basis vectors.)
\end{proof}

% ═════════════════════════════════════════════════════════════════════════════
\section{Single-Bit Error Correction}
% ═════════════════════════════════════════════════════════════════════════════

\begin{proposition}
$C$ is single-bit error correcting, i.e.\ the minimum Hamming distance satisfies
$d_{\min} \geq 3$.
\end{proposition}

\begin{proof}
For a binary linear code with parity check matrix $A$, $d_{\min} \geq 3$ if and
only if every column of $A$ is nonzero and no two columns of $A$ are equal.
The six columns of $A$ are
\[
    \mathbf{a}_1 = \begin{pmatrix}1\\1\\1\\0\end{pmatrix},\quad
    \mathbf{a}_2 = \begin{pmatrix}0\\1\\1\\0\end{pmatrix},\quad
    \mathbf{a}_3 = \begin{pmatrix}1\\1\\0\\1\end{pmatrix},\quad
    \mathbf{a}_4 = \begin{pmatrix}0\\0\\1\\1\end{pmatrix},\quad
    \mathbf{a}_5 = \begin{pmatrix}0\\1\\0\\1\end{pmatrix},\quad
    \mathbf{a}_6 = \begin{pmatrix}1\\0\\0\\0\end{pmatrix}.
\]
All six are nonzero, and inspection shows they are pairwise distinct. Hence
$d_{\min} \geq 3$, and $C$ can correct any single-bit transmission error.
\end{proof}

% ═════════════════════════════════════════════════════════════════════════════
\section{Syndrome Decoding}
% ═════════════════════════════════════════════════════════════════════════════

When a single-bit error occurs in position $i$, the received word $\mathbf{r}$
satisfies $A\mathbf{r} = \mathbf{a}_i$ (the $i$-th column of $A$), because
\[
    A(\mathbf{c} + \mathbf{e}_i) = A\mathbf{c} + A\mathbf{e}_i
    = \mathbf{0} + \mathbf{a}_i = \mathbf{a}_i,
\]
where $\mathbf{c}$ is the sent codeword and $\mathbf{e}_i$ is the standard
basis vector with a $1$ in position $i$. The syndrome $\mathbf{s} = A\mathbf{r}$
identifies the error position by matching $\mathbf{s}$ to the columns of $A$.

\subsection{Syndrome lookup table}

\begin{center}
\renewcommand{\arraystretch}{1.2}
\begin{tabular}{cc}
    \hline
    \textbf{Syndrome} $\mathbf{s}$ & \textbf{Error position} \\
    \hline
    $(1,1,1,0)^{\top}$ & 1 \\
    $(0,1,1,0)^{\top}$ & 2 \\
    $(1,1,0,1)^{\top}$ & 3 \\
    $(0,0,1,1)^{\top}$ & 4 \\
    $(0,1,0,1)^{\top}$ & 5 \\
    $(1,0,0,0)^{\top}$ & 6 \\
    \hline
\end{tabular}
\end{center}

\subsection{Decoding received words}

\begin{example}
For each received word below, determine the transmitted codeword:
\[
    \mathbf{a} = (1,1,1,1,1,1)^{\top}, \quad
    \mathbf{b} = (0,1,0,1,0,1)^{\top}, \quad
    \mathbf{c} = (0,0,0,1,1,1)^{\top}.
\]
\end{example}

\begin{proof}[Solution]
We compute the syndrome for each received word and consult the lookup table.

\textbf{Word $\mathbf{a} = (1,1,1,1,1,1)^{\top}$:}
\[
    A\mathbf{a} \equiv
    \begin{pmatrix}
        1{+}0{+}1{+}0{+}0{+}1\\
        1{+}1{+}1{+}0{+}1{+}0\\
        1{+}1{+}0{+}1{+}0{+}0\\
        0{+}0{+}1{+}1{+}1{+}0
    \end{pmatrix}
    \equiv
    \begin{pmatrix}1\\0\\1\\1\end{pmatrix} \pmod{2}.
\]
The syndrome $(1,0,1,1)^{\top}$ does not match any entry in the lookup table,
so $\mathbf{a}$ does not have a single-bit error (multiple errors have occurred
or transmission was uncorrectable). No unique correction exists.

\medskip
\textbf{Word $\mathbf{b} = (0,1,0,1,0,1)^{\top}$:}
\[
    A\mathbf{b} \equiv
    \begin{pmatrix}
        0{+}0{+}0{+}0{+}0{+}1\\
        0{+}1{+}0{+}0{+}0{+}0\\
        0{+}1{+}0{+}1{+}0{+}0\\
        0{+}0{+}0{+}1{+}0{+}0
    \end{pmatrix}
    \equiv
    \begin{pmatrix}1\\1\\0\\1\end{pmatrix} \pmod{2}.
\]
The syndrome $(1,1,0,1)^{\top} = \mathbf{a}_3$ indicates an error in position $3$.
Flipping bit 3: $\mathbf{b}' = (0,1,1,1,0,1)^{\top}$ is the transmitted codeword.
We verify: $A\mathbf{b}' \equiv (0,0,0,0)^{\top}$.

\medskip
\textbf{Word $\mathbf{c} = (0,0,0,1,1,1)^{\top}$:}
\[
    A\mathbf{c} \equiv
    \begin{pmatrix}
        0{+}0{+}0{+}0{+}0{+}1\\
        0{+}0{+}0{+}0{+}1{+}0\\
        0{+}0{+}0{+}1{+}0{+}0\\
        0{+}0{+}0{+}1{+}1{+}0
    \end{pmatrix}
    \equiv
    \begin{pmatrix}1\\1\\1\\0\end{pmatrix} \pmod{2}.
\]
The syndrome $(1,1,1,0)^{\top} = \mathbf{a}_1$ indicates an error in position $1$.
Flipping bit 1: $\mathbf{c}' = (1,0,0,1,1,1)^{\top}$ is the transmitted codeword.
We verify: $A\mathbf{c}' \equiv (0,0,0,0)^{\top}$.
\end{proof}
